\chapter{From ANCIENT to MODERN Times}\label{ch:ancient}
\markboth{From Ancient to Modern Times}{From Ancient to Modern Times}

\chapquote{``We are made of stardust, our whole body consists of material that has been here before the beginning of time.''}{Giorgio A.~Tsoukalos}

% ---------------------------------------------------------------
\section{Classifying Impossibilities}\label{sec:classifyingimpossibiliti}

Graham Hancock has spent thirty years making one argument, and it deserves to be stated properly
before anybody argues with it. The claim is not that aliens built the pyramids. It is that a
sophisticated human civilisation existed before the end of the last Ice Age, that it was destroyed
in the cataclysms of the Younger Dryas around 12{,}900 years ago, and that what survived was
carried by refugees into the cultures we call the beginnings of history --- as astronomy encoded in
myth, as agriculture appearing suspiciously fully formed, as megalithic building traditions
separated by oceans. On this reading, the ancient world is not the dawn of human achievement. It
is the amnesia after a collapse.

It is a genuinely interesting thesis and it fails on evidence rather than on principle, which is a
compliment. A civilisation leaves middens, quarries, harbours, pollen signatures, metallurgical
slag, genetic bottlenecks, isotopes in ice. We have found none of that from the right period, and
we have looked at that period intensively for other reasons. G\"obekli Tepe --- Hancock's best
exhibit, and it is a good one --- is around 11{,}500 years old, monumental, sophisticated, and
built by people whose toolkit and diet we can characterise in detail. It pushed the date of
monumental architecture back by millennia. It did not produce a lost Atlantis; it produced
hunter-gatherers doing something nobody expected hunter-gatherers to do. Chapter~\ref{ch:atlantis} returns to
the submerged-city version of the same idea.

I open with Hancock deliberately, because his structure is the structure of this entire chapter.
Something in the ancient record looks impossible. A modern reader supplies the missing cause. And
whichever cause you supply --- lost civilisation, visiting spacecraft, divine intervention --- you
are performing the same operation. The useful skill is not picking a favourite. It is sorting the
impossibilities before explaining them: which are impossible because the evidence is
genuinely anomalous, which are impossible only relative to what the ancient author knew, and which
are impossible only because somebody described them badly a very long time ago and the description
is all we have left.

Begin with the middle category, because it comes with a ready-made test.

A miracle, in the strict theological sense, is an event caused by direct divine intervention in
violation of natural law. That definition is precise and it has a consequence people rarely
follow through on: a miracle is only a miracle relative to what the observer knows about natural
law. Move the observer, and the category moves with them.

The Ancient Alien Hypothesis supplies spacecraft. Religion supplies God. Hancock supplies a
drowned civilisation. Conventional history supplies exaggeration, metaphor, and craftsmanship we
have underestimated. All four are the same move --- filling an explanatory gap with the most
impressive thing available to the person doing the filling. Only one of those four is required to
suspend a law of nature, and there is a classical test for that case.

So before we go further, the ground rules from Chapter~\ref{ch:language}, applied to text:

Ancient documents are not neutral reports. They are written by literate elites, for purposes ---
legitimising a dynasty, glorifying a deity, explaining a ruin nobody remembered building.
They are copied by hand for centuries, with drift. They are translated, often several times,
and translation of technical vocabulary from a dead language is a guessing game.

\hi{And the correct null hypothesis for an ancient account of something amazing is not ``it didn't
happen''.} It is: something happened, and the description we have is the last surviving link in a
long chain of retelling. Ezekiel saw something, or believed he did, or wrote a vision for a
reason. Determining which is history; determining what he \emph{physically} saw is, in almost
all cases, not recoverable.

\begin{funfact}
David Hume gave the sharpest test in 1748, and it is still the one to use. The question is never
``is this event possible?'' but ``which is more likely --- that the laws of nature were
suspended, or that the testimony is mistaken?'' Human testimony is demonstrably fallible; the
laws of nature have never once been observed to fail. This is not an argument that miracles
cannot happen. It is an argument about where to place your bet before the evidence arrives.
\end{funfact}

% ---------------------------------------------------------------
\section{The Ancient Alien Hypothesis}\label{sec:aah}

The Ancient Alien Hypothesis holds that extraterrestrials visited Earth in antiquity, were
recorded as gods, and contributed technology or genetic material to early humans. Its modern
popular form comes from Erich von D\"aniken's \emph{Chariots of the Gods?} (1968), which sold
something like 70 million copies, though the idea predates him --- Charles Fort gestured at it,
and Morris Jessup and the Soviet writer Matest Agrest published versions in the 1950s.

Let me first state the strongest version of the case, because there is one and it is worth
engaging with.

Sky gods are close to universal. Cultures with no contact describe beings descending from above,
teaching agriculture, writing, metallurgy and law, and departing with a promise to return. Some
ancient engineering is genuinely difficult to account for --- the largest stones at Baalbek
weigh over 800 tonnes and one unmoved block exceeds 1,\!600. \hi{Some ancient astronomical
knowledge is precise beyond apparent need.} And the Dogon of Mali were reported by Marcel
Griaule to possess knowledge of Sirius B, a white dwarf invisible without a telescope.

Now the problems, in order of severity.

\textbf{It is an argument from ignorance, structurally.} Every instance takes the form: I cannot
see how they did this, therefore someone else did. \hi{That inference is invalid regardless of the
candidate you insert.} It also has a poor track record, because the gaps keep closing --- and
Section~\ref{sec:impossiblestone} is about exactly that.

\textbf{It underestimates ancient people, systematically and in one direction.} The
hypothesis has never once been applied to European megaliths in a form its adherents find
equally exciting, and it is applied overwhelmingly to non-European achievement: Egyptian,
Mesoamerican, Andean, Southeast Asian. There is a name for the assumption that brown people
could not have built their own monuments, and the fact that its modern proponents do not intend
it does not make the pattern less real. Egyptians left us quarry marks, workers' villages,
bakeries, medical records of construction injuries, administrative papyri recording the
delivery of limestone blocks --- the Merer diary, found at Wadi al-Jarf in 2013, is a working
logbook of a boat crew hauling casing stone to Giza. We know who built the pyramids. Their
names are on the walls, and some of them scratched graffiti about which work gang was better.

\textbf{The material evidence is absent, and it is the kind of absence that counts.} A visiting
technological civilisation would leave residue that survives: alloys that do not occur
naturally, isotopic anomalies, machined tolerances, plastics, radioactive traces, refined
materials in stratified archaeological context. We have excavated an enormous number of sites
with increasingly sensitive methods and found none of it. Not one anomalous artefact from a
controlled excavation --- as Section~\ref{sec:deconstructingpeerreview} explains, that is where such a find would have to come
from to mean anything.

\textbf{The Dogon case dissolved.} Later anthropologists, notably Walter van Beek, worked with
the Dogon and could not find the Sirius B tradition Griaule reported. Griaule had visited in
the 1930s and 1940s --- after Sirius B was well publicised in European popular science, and
after European visitors including astronomers had been in the region. The likeliest explanation
is feedback from Griaule's own questioning, which is Section~\ref{sec:confirmationbias}'s leading-question problem in
fieldwork rather than in a therapist's office.

\subsection{Sitchin, the Anunnaki, and the Twelfth Planet}

If von D\"aniken supplied the hypothesis, Zecharia Sitchin supplied its scripture, and since his
names turn up constantly in this field you should know exactly where they come from. Sitchin was a
Russian-born American writer who published \emph{The 12th Planet} in 1976 and eventually seven
volumes of what he called the Earth Chronicles. His claim was that the Sumerian cuneiform record
--- the oldest written literature we possess, from the oldest civilisation we can name --- is not
mythology but history. The Anunnaki, in his reading, were flesh-and-blood visitors from Nibiru, a
twelfth planet on a 3,\!600-year elliptical orbit, who came to Earth to mine gold and
genetically engineered humanity from local hominids as a labour force. Every element of the modern
ancient-astronaut vocabulary --- Anunnaki, Nibiru, the gold mining, the engineered slave species
--- is Sitchin's invention, not Sumer's, and he was a fixture on \emph{Coast to Coast AM} under
both Art Bell and George Noory, which is how it reached millions.

\figwrap[l]{0.38}{A Sumerian administrative tablet of the kind Sitchin cited. The signs he read as ``rocketship'' say nothing of the sort.}{https://commons.wikimedia.org/wiki/File:Sumerian_cuneiform_tablet.jpg}{fig:cuneiform}
The trouble with Sitchin is not the one usually raised.

\actualq[problem]{that Sitchin was a kook, though mainstream Assyriology regarded him as one.}{that
the translations are wrong --- specifically, checkably wrong.}

Sitchin was self-taught in the cuneiform languages (see Figure~\ref{fig:cuneiform}) and worked without the philological apparatus. Assyriologists ---
Michael Heiser wrote the most detailed rebuttal --- have gone through his key terms line by line.
The word he renders as ``rocketship'' does not mean that. \emph{Anunnaki} means, roughly,
those of princely seed, and the Sumerian texts place them in the sky and the underworld as deities,
not in vehicles. There is no Sumerian word for Nibiru's supposed orbit, no tablet listing twelve
planets, and the Sumerians did not know Neptune or Uranus existed. When a body of work rests on
translation and the translation fails, nothing downstream of it survives. That is not a
sceptic's opinion; it is the ordinary standard applied to any scholar who works from a text.

As for Nibiru, the ufological literature likes to say NASA supports the idea of a large
gravitational body beyond the Kuiper belt. What NASA-affiliated astronomers have actually
proposed --- Batygin and Brown, in 2016 --- is Planet Nine: a hypothetical body of perhaps five to
ten Earth masses, inferred from the clustered orbits of distant trans-Neptunian objects, at
several hundred astronomical units, with a period on the order of ten thousand years or more. It
has never been observed, its existence is contested, and even if confirmed it would be a frozen
world with no possible relation to Sumerian labour disputes. A distant, cold, unconfirmed planet
and an inhabited world that swings through the inner solar system every 3,\!600 years are not the
same claim wearing different clothes. This is a move you will see again and again, so learn it
now: take a real, cautious, technical result, keep the headline, and quietly swap the content.

The same surgery is worth performing on Robert Bauval, whose \emph{Orion Correlation Theory} argues
that the three Giza pyramids map the belt stars of Orion. The correlation is loose, requires
choosing a particular epoch and ignoring the pyramids that do not fit, and the Egyptians left us
their own extensive account of what the alignments were for --- the northern circumpolar stars and
the king's transformation --- which is both better attested and more interesting. Bauval's water
erosion argument about the Sphinx enclosure, developed with Robert Schoch, is the most respectable
thing in this territory and the one place where a real geological dispute exists. It is also not an
argument about aliens, and no Egyptologist has yet been persuaded that it beats the conventional
dating.

\begin{funfact}
The sky-god universality argument has a much better explanation than visitors, and it is the
same one that explains why Section~\ref{sec:geneticmotivation}'s abduction narrative matches fairy folklore. Every human
culture is under the same sky, subject to the same weather, dependent on the same sun,
frightened by the same lightning, and possessed of the same brain --- one that models the world
in terms of agents with intentions, because that heuristic kept our ancestors alive. Universal
myths are evidence of universal humans. That is a real discovery about our species, and it is
more interesting than a spaceship.
\end{funfact}

% ---------------------------------------------------------------
\section{UFOs in ANCIENT Artwork}\label{sec:ufosancientartwork}

The standard gallery is short and every item in it has an art-historical answer. I will go
through them because knowing the answers is the only inoculation.

\textbf{\emph{The Madonna with Saint Giovannino}} (15th century, Palazzo Vecchio). A disc-shaped
object with rays in the sky, a man and his dog looking up. It is a standard Renaissance
depiction of a cloud with divine radiance --- the visual convention for the Annunciation or the
glory of God --- and comparable clouds appear in dozens of contemporary works nobody has heard
of. The shepherd and dog are conventional too.

\textbf{\emph{The Baptism of Christ}} by Aert de Gelder (1710). A bright disc emitting beams
over Christ's head. This is the Holy Spirit descending, which the text of Matthew 3:16 requires
the painter to depict. The alternative is that de Gelder painted a UFO into a baptism scene and
nobody in 1710 remarked on it.

\textbf{\emph{The Crucifixion}} (1350, Visoki De\v{c}ani monastery, Kosovo). Two objects with
figures apparently inside them, flanking the cross. They are personifications of the Sun and
Moon --- a standard Byzantine convention with a long documented pedigree, present in countless
Orthodox crucifixion scenes, representing creation mourning.

\textbf{The Nazca lines.} Enormous geoglyphs in the Peruvian desert, only fully legible from the
air, therefore runways or signals. They are not runways: they are shallow surface scrapings a
few centimetres deep in loose stone, which would not survive a bicycle let alone a landing.
They are constructible from the ground with stakes and cord, which has been demonstrated
repeatedly, and are visible from surrounding foothills. They are almost certainly ritual
pathways related to water and were walked, not viewed.

\textbf{The Palenque sarcophagus lid} (Maya, 7th century). K'inich Janaab' Pakal reclining,
read as an astronaut in a capsule. The full iconography is documented and internally consistent:
Pakal falls into the underworld at the moment of death, with the world tree rising above him,
the Celestial Bird at its top, and the maw of the underworld below. The glyphs around the border
name him and date it. Rotate the image ninety degrees to make it look like a rocket and you have
rotated a religious scene out of its own orientation.

\textbf{Ezekiel's wheel.} Ezekiel 1 describes wheels within wheels, full of eyes, with four
living creatures. Josef Blumrich, a NASA engineer, published \emph{The Spaceships of Ezekiel} in
1974 reconstructing a lander from the text. The text is apocalyptic vision literature written in
Babylonian exile, deploying Babylonian and Assyrian imagery --- the four-faced creatures
correspond closely to known Mesopotamian composite guardian figures --- and its purpose is
theological: God's throne is mobile, therefore God is present in exile, outside the destroyed
Temple. That is the point of the passage and it is a devastating one for its original audience.
Section~\ref{sec:ezekiel} takes Blumrich's reconstruction apart properly, because it is the best
version of this whole category and deserves the space.

\textbf{The Dendera light} and the ``helicopter hieroglyphs'' at Abydos. The former is a lotus
and a serpent, a standard creation motif, in a crypt whose full inscriptions explain it. The
latter is a genuine curiosity with a dull answer: it is a palimpsest, where Ramesses II's titles
were carved over Seti I's and the filling plaster fell out, superimposing two sets of glyphs.
This is visible on inspection and documented in the epigraphic record.

\begin{funfact}
Here is the test that settles the whole category. Ancient artists depicted what they saw with
remarkable fidelity when they were depicting objects --- Egyptian tomb painting gets bird
species right to the feather, and Greek pottery renders ship rigging accurately enough to
reconstruct from. If flying craft had been observed, we would have detailed, repeated,
unambiguous depictions with consistent features, in secular as well as religious contexts, and
they would be described in the accompanying inscriptions. Instead we have a handful of religious
symbols that require reinterpretation against their own captions.
\end{funfact}

% ---------------------------------------------------------------
\section{Ley Lines}\label{sec:leylines}

In 1921, an English businessman and amateur photographer named Alfred Watkins was looking at
a map of Herefordshire when he noticed that certain ancient sites --- standing stones,
mounds, churches on older sacred sites, hilltop beacons --- appeared to fall into straight
lines. He published \emph{The Old Straight Track} in 1925, proposing that these alignments
were prehistoric trading routes surveyed by sighting between landmarks.

Watkins' original hypothesis was archaeological and entirely reasonable in outline. People do
navigate by sightlines. Some ancient landscapes do contain deliberate alignments --- the
solstice axis at Stonehenge and Newgrange are real, measured, and beyond dispute.

The problem is statistical, and it is a lovely example of why Chapter~\ref{ch:language} exists. England
contains many thousands of ancient sites. If you scatter that many points randomly on a map
and permit a line to be a few hundred metres wide, you will find an enormous number of
``alignments'' by pure chance. This was demonstrated formally --- in 1980 the astronomer
Robert Forrest ran the numbers and showed that the observed frequency of Watkins-style
alignments is what you would expect from random placement, and in 1989 the mathematician
Chris Marshall generated convincing ley networks from randomly generated site data. You can
do this yourself with a map of pizza restaurants, and several people have, producing
spectacular ``ley lines'' across Manhattan.

By the 1960s the concept had been detached from archaeology entirely and reattached to earth
energies, dowsing, and eventually to UFOs --- the claim being that sightings cluster along
leys, that craft navigate by them, and that they represent lines of geomagnetic force. There
is no measured geomagnetic anomaly along any ley line. Ordnance Survey and geological survey
data have been available for a century, and the alignments do not correspond to anything in
them.

\begin{funfact}
The UFO-ley connection has an author: Tony Wedd, in 1961, who combined Watkins with the
French ``orthot\'enie'' work of Aim\'e Michel, who had claimed in 1958 that sightings during
the 1954 French wave fell along straight lines. Michel's orthoteny was tested properly in
the 1960s by Jacques Vall\'ee among others, and it collapsed: the lines were an artefact of
selecting which reports to plot. Vall\'ee, to his considerable credit, published the negative
result about his own earlier enthusiasm. That is what Section~\ref{sec:investigativetechniques} step six looks like in
practice.
\end{funfact}

I keep ley lines in the book for one reason: they are the cleanest demonstration in the
entire field that pattern-finding is not evidence. \hi{Your brain is extremely good at spotting
alignments and extremely bad at estimating how many alignments chance produces.} Every time
you see a map in this literature with lines drawn on it --- and you will see many, connecting
pyramids, connecting crash sites, connecting sacred mountains --- ask the only question that
matters: how many lines would I have found in random data? Nobody who draws those maps has
ever once told you.

% ---------------------------------------------------------------
\section{The Impossible Stone}\label{sec:impossiblestone}

The strongest argument the ancient astronaut literature has is not artwork. It is masonry, and
I want to give it a fair hearing, because standing in front of some of these walls does
something to your confidence that reading about them does not.

\figwrap[r]{0.42}{The Trilithon at Baalbek: three courses of roughly 800 tonnes each, moved by documented Roman engineering.}{https://commons.wikimedia.org/wiki/File:Baalbek_trilithon.jpg}{fig:baalbek}
The items usually cited: the Great Pyramid, 2.3 million blocks, some casing stones fitted to
tolerances measured in fractions of a millimetre. The Trilithon at Baalbek (see Figure~\ref{fig:baalbek}), three blocks around
800 tonnes each, raised into a wall course, with the abandoned Stone of the Pregnant Woman and
a larger 2014 find nearby exceeding 1,\!600 tonnes. Sacsayhuam\'an, polygonal Inca walls of
irregular multi-tonne blocks fitted without mortar so closely that a knife blade will not
enter. The Osirion at Abydos. The unfinished obelisk at Aswan, which would have weighed around
1,\!100 tonnes. Puma Punku's andesite H-blocks, with flat faces and interior right angles.

The conventional answers are unglamorous, well tested, and --- this is the part people miss ---
experimentally demonstrated.

\textbf{Cutting.} Egyptian hard-stone work used copper saws with quartz sand abrasive, and it
is the sand doing the cutting, not the metal. Drill cores from Egyptian granite show the
characteristic striations. Denys Stocks spent decades replicating Egyptian stoneworking with
period tools and successfully cut, drilled and dressed granite. Andesite at Puma Punku was
worked with harder hammerstones and abrasive; the ``machined'' interior angles show tool marks
under examination and are not, in fact, perfectly square.

\textbf{Fitting.} Inca polygonal masonry was achieved iteratively --- place a block, mark the
contact, lift, dress, replace, repeat. It is laborious rather than mysterious, and unfinished
Inca walls survive showing the process mid-sequence, with dressing scars and abandoned blocks.
The technique also has a functional reason nobody mentions: irregular interlocking mortarless
masonry is superbly earthquake resistant, and it is still standing in a seismic zone that has
flattened everything the Spanish built.

\textbf{Moving.} This is where the intuition fails hardest, and it fails because of
arithmetic. A 60-tonne block on a lubricated wooden sledge on a prepared trackway requires a
pulling force well within the reach of a few hundred coordinated men --- a wall painting in the
tomb of Djehutihotep depicts 172 men hauling an alabaster colossus, with a man pouring liquid
in front of the sledge, and a 2014 experiment at Amsterdam confirmed that wetting sand roughly
halves the required force. For the Baalbek trilithon, the Roman engineering literature and
surviving evidence of capstans, levers, ramps and earth embankments account for it; the blocks
were quarried a few hundred metres uphill from their final position, which is a hint the
``impossible'' framing always omits.

\textbf{Precision.} Ancient surveying was extremely good. The Great Pyramid's base is level to
around two centimetres across 230 metres, which is achievable with water channels, and its sides
are aligned to true north within about three arcminutes, achievable by stellar observation.
These are not exotic methods. They are patient ones.

\begin{funfact}
What these sites really demonstrate is organisation, not technology --- and organisation is the
harder achievement. The Great Pyramid required feeding, housing, sanitising and administering a
workforce of tens of thousands for two decades, with supply lines, bakeries, breweries,
medical care, and a bureaucracy tracking it all in writing. We found the workers' town at Giza,
with its bakeries and its cattle bones and the fractures set by physicians. The astonishing
thing about the pyramids is not that stone can be cut. It is that human beings four and a half
thousand years ago could run a project of that scale at all. Handing that to visitors is not
generous to the ancients. It is theft.
\end{funfact}

% ---------------------------------------------------------------
\section{The Impossible Glass}\label{sec:impossibleglass}

There is a class of claim about desert glass that gets cited as evidence of ancient nuclear
warfare, and it is worth taking apart because it involves real physical evidence with a real
explanation --- my favourite kind of case.

\textbf{Libyan Desert Glass.} Sheets and fragments of nearly pure silica glass scattered across
about 6,\!500 square kilometres of the Great Sand Sea between Libya and Egypt. It is roughly 98\%
silica, yellow-green, and about 29 million years old. Tutankhamun's pectoral contains a carved
scarab of it, which is a genuinely wonderful fact. Its origin is impact-related: it contains
iridium enrichment, high-temperature mineral phases including cristobalite, and traces of
meteoritic material, indicating formation at temperatures above 1,\!600$^\circ$C, probably from
an airburst or a nearby impact melt. Twenty-nine million years predates humans by a wide
margin, which disposes of the ancient war reading without further argument.

\textbf{Trinitite}, by contrast, is genuine nuclear glass --- formed at the Trinity test site in
July 1945, greenish, and radioactive, containing fission products and traces of the device and
its tower. Comparing the two is instructive. Nuclear glass carries an unmistakable signature:
caesium-137, strontium-90, plutonium isotopes, europium-152 from neutron activation. It is
detectable for millennia and it is not subtle.

So when Mohenjo-daro is claimed as an ancient nuclear site --- vitrified surfaces, skeletons in
the streets, elevated radioactivity --- the test is available and has been applied. There is no
elevated radioactivity at Mohenjo-daro. The claim traces to a single unsourced assertion
repeated for decades without anybody measuring. The site shows evidence of decline consistent
with hydrological change in the Indus system, and the ``skeletons in the street'' are burials
and deposits from different strata that were reported carelessly in early excavation.

\textbf{Vitrified forts} in Scotland and elsewhere are real: Iron Age dry-stone ramparts whose
stones have partially melted and fused. Experimental archaeology has reproduced the effect by
burning timber-laced ramparts, and the debate in the literature is whether the burning was
destruction by enemies, ritual, or deliberate construction technique --- a real scholarly
disagreement about temperature and duration, being argued with reproducible fires.

\textbf{Fulgurites} --- lightning-strike glass --- form tubes in sand where a strike has fused
silica along its path, and account for a great deal of ``anomalous'' glass in deserts.

\begin{funfact}
The Mahabharata's descriptions of the Brahmastra --- ``a single projectile charged with all the
power of the Universe... an incandescent column of smoke and flame as bright as ten thousand
suns'' --- are genuinely striking, and often quoted. The quotation is also, in most circulating
versions, not a translation of the Mahabharata. It is a paraphrase that entered the literature
through mid-twentieth-century popular writing, after Hiroshima, by authors who had a modern
image in their heads. Compare it to a scholarly translation and the resemblance largely
evaporates. Oppenheimer did quote the Gita at Trinity --- ``Now I am become Death'' --- and the
feedback loop starts there.
\end{funfact}

% ---------------------------------------------------------------
\section{The Strange Case of Brawley Oates}\label{sec:strangecasebrawleyoates}

In the early 1960s, in Elberton, Georgia, a man named Brawley Oates bought the property
surrounding a small granite outcrop called Wonder Hill --- a spot where, according to local
legend and to a good many visitors, things do not behave as they should. Water appears to run
uphill. A ball rolls the wrong way. A person standing in the structure leans at an impossible
angle. Oates built the attraction that later generations knew as a ``gravity hill'' house, and
the story attached to him afterwards is that he drank the water from the site to cure his
arthritis and that it worked.

I include this case for a reason that has nothing to do with whether Oates was cured. Gravity
hills --- and there are hundreds catalogued worldwide, from the Electric Brae in Ayrshire to
Magnetic Hill in New Brunswick to Confusion Hill in California --- are the single best available
demonstration that direct sensory experience is unreliable in a completely specific and
measurable way.

\figwrap[l]{0.38}{A gravity hill with the true horizon obscured: the road that looks uphill surveys as downhill.}{https://commons.wikimedia.org/wiki/File:Gravity_hill_road.jpg}{fig:gravityhill}
Here is what is happening. Human beings do not sense the gravitational vertical directly with
any precision. We infer level from visual cues: the horizon, the angle of trees, the line of a
road, the walls of the room we are standing in. Obstruct the true horizon --- with a slope, a
ridgeline, a stand of trees growing at an angle on a hillside, or a deliberately tilted
building --- and the visual system will confidently construct a false reference frame, and
your body will believe it. Surveying instruments at gravity hills (see Figure~\ref{fig:gravityhill}) consistently show that the
apparently uphill slope is in fact downhill. The water is obeying physics. Your inner ear and
your visual cortex are having an argument, and vision wins.

The purpose-built versions --- ``mystery spots'' --- amplify this with rooms constructed out of
plumb, so that every reference line in the observer's field of view is tilted consistently.
Stand in one and you will experience an overwhelming sensation of leaning, and you will not be
able to reason your way out of it. That is the whole point of visiting.

\begin{funfact}
Elberton, Georgia has a genuinely stranger claim to fame that requires no anomaly at all. In
1980 a man using the name R.C.~Christian commissioned, and paid for, an enormous granite
monument outside town --- the Georgia Guidestones --- inscribed in eight languages with ten
instructions for rebuilding civilisation after a collapse, including maintaining humanity under
500 million. Nobody ever established who paid for it. It became a fixture of
New World Order literature for four decades, and in July 2022 somebody bombed it and the
remains were demolished the same day. A real unsolved mystery, in the same small town, with
better documentation than any UFO case in this book.
\end{funfact}

The methodological point stands on its own: if a hillside in Georgia can convince every visitor
that water flows uphill, then ``I saw it with my own eyes'' establishes considerably less than
people think. This is Section~\ref{sec:correlationvscausation}'s argument, standing outdoors, available for anybody to test
with a spirit level.
